\documentclass{article}
\usepackage{iclr2026_conference,times}
\input{math_commands.tex}

\usepackage{amsmath}
\usepackage{amssymb}
\usepackage{booktabs}
\usepackage{multirow}
\usepackage{graphicx}
\usepackage{microtype}
\usepackage{xcolor}
\usepackage{hyperref}
\usepackage{url}

\title{From Channel State Information to Language:\\
Physics-Grounded Multi-Task Learning for Interpretable Wireless Channel Descriptions}

\author{Anonymous Authors}

% Keep anonymous for submission.
% \iclrfinalcopy

\newcommand{\model}{\textsc{CSI2Text}}
\newcommand{\mae}{\operatorname{MAE}}
\newcommand{\los}{\textsc{LoS}}
\newcommand{\nlos}{\textsc{NLoS}}
\newcommand{\etal}{et al.\ }

\begin{document}
\maketitle

\begin{abstract}
Channel state information (CSI) is information-rich but difficult for a human to interpret directly.
We study the conversion of high-dimensional CSI into concise natural-language descriptions grounded in propagation physics.
Our approach, \model, combines a configuration-aware Transformer encoder, contrastive alignment between CSI and structured channel descriptions, a shared physics representation, and specialized prediction branches for delay, angle, power, $K$-factor, and reflection statistics.
The predicted attributes are verbalized through a deterministic, auditable template, followed by a lightweight relational correction that removes contradictions such as a first-path delay preceding the direct line-of-sight delay.
On a balanced 100k-sample UPA-$8{\times}8$, 128-subcarrier dataset, the full multi-task model obtains a first-path delay MAE of $21.91\pm0.30$ ns, first-path angle MAE of $18.37\pm0.37^\circ$, and reflection-interaction exact accuracy of $70.61\pm0.39\%$ over three seeds.
On a separately constructed, naturally imbalanced reflected-path extension, the three-seed generated descriptions reach $94.06\pm1.07\%$ composite factuality, $99.97\pm0.01\%$ numeric-slot F1, $0.04\pm0.02\%$ numeric-slot hallucination, and $99.87\pm0.03\%$ physical consistency.
The extension predicts the number of reflected paths with a text-level MAE of $0.091\pm0.007$.
These results indicate that structured language can serve as a faithful and inspectable interface to learned wireless channel representations.
\end{abstract}

\section{Introduction}

CSI describes how a wireless channel transforms a transmitted signal across antennas and frequencies.
It supports beam management, localization, sensing, and channel modeling, yet its complex-valued tensor form is inaccessible to most users and difficult to audit.
Existing learning systems typically map CSI to a task label or a small set of regression targets.
Although useful, such outputs do not provide a unified, human-readable account of the propagation scene.

We instead ask: \emph{can a model translate CSI into a factual language description of the channel?}
A desired output may state that a signal is indoor and \nlos, followed by the estimated first-path delay and angle, $K$-factor, delay and angular spreads, path richness, and reflection statistics.
This formulation differs from unconstrained captioning.
Wireless attributes have units, conditional semantics, and known relations; fluent text is not sufficient if the stated values contradict one another or the reference channel.

We introduce \model, a physics-grounded multi-task framework that learns a common representation for CSI, language prototypes, and propagation attributes.
The CSI encoder operates on beam-frequency measurements and conditions its tokens on acquisition configuration.
Three contrastive objectives align CSI, instance descriptions, and semantic prototypes.
Specialized heads then model heterogeneous targets using appropriate representations: delay-profile features for temporal quantities, power-profile features for received power, and circular sine/cosine targets for angles.
Finally, an explicit verbalizer renders a structured record into natural language.
Because every phrase can be traced to a predicted slot, the output remains auditable.

Our contributions are:
\begin{itemize}
    \item We formulate CSI-to-language conversion as grounded structured prediction followed by faithful verbalization, rather than unconstrained text generation.
    \item We develop a configuration-aware, physics-guided multi-task architecture with language alignment, shared physical features, and target-specific branches.
    \item We define a slot-based factuality protocol that separately measures categorical correctness, slot coverage, tolerance-aware numerical accuracy, hallucination, and physical consistency.
    \item We provide three-seed comparisons, ablations, a relational-correction analysis, and a reflected-path-count extension while explicitly separating the balanced main dataset from the imbalanced extension dataset.
\end{itemize}

\section{Related Work}

\paragraph{Learning from CSI.}
Public ray-tracing resources such as DeepMIMO \citep{alkhateeb2019deepmimo} and RayVerse \citep{rayverse2025} facilitate reproducible learning from channel measurements with path-level annotations.
Most CSI learning pipelines optimize a single downstream task, whereas our goal is to expose several coupled propagation properties through one readable description.
Physics-informed wireless modeling has recently emphasized combining physical structure with learned representations to improve interpretability and generalization \citep{zhu2024physicsinformed}.

\paragraph{Cross-modal and language-aligned representation learning.}
Contrastive language--signal alignment follows the general principle that paired modalities can share a semantic embedding space, as demonstrated at scale for images and text by CLIP \citep{radford2021clip}.
Our setting differs because CSI has no naturally authored captions and the language is derived from measurable propagation attributes.
We therefore align CSI not only with instance descriptions but also with semantic prototypes and supervise the latent space using explicit physics targets.

\paragraph{Wireless signals and language.}
Recent work has connected radio-frequency measurements to language for human-centric semantic tasks.
WiFi2Cap produces action captions from Wi-Fi CSI \citep{wei2026wifi2cap}, WirelessSenseLLM aligns wireless sensing with language models for activity understanding \citep{keya2026wirelesssensellm}, and RF-GPT maps RF observations to structured semantic outputs \citep{zou2026rfgpt}.
In contrast, \model focuses on the \emph{propagation channel itself}: the generated description is a factual rendering of channel geometry and statistics rather than an open-ended description of human behavior.

\section{Method}
\label{sec:method}

\subsection{Problem formulation}

Let $\mathbf{H}\in\mathbb{C}^{B\times F}$ denote beam-domain CSI with $B$ beam tokens and $F$ frequency samples, and let $\mathbf{c}$ contain acquisition metadata such as bandwidth and frequency configuration.
Each example has a structured physical record
\begin{equation}
  \mathbf{y} = (e,\ell,n_p,\tau_1,\theta_1,p_1,K,\sigma_\tau,\sigma_\theta,
  \tau_{\mathrm{LoS}},\theta_{\mathrm{LoS}},n_r,n_{rp}),
\end{equation}
where $e$ is environment, $\ell$ is \los/\nlos\ status, $n_p$ is path count, $\tau_1,\theta_1,p_1$ are first-path delay, angle, and power, $K$ is the Rician $K$-factor, $\sigma_\tau,\sigma_\theta$ are delay and angular spreads, $n_r$ is the \emph{total number of reflection interactions}, and $n_{rp}$ is the \emph{number of paths containing at least one reflection}.
The last quantity is used only in the reflected-path extension.
This distinction is important: a single path may undergo several reflections, so $n_r\le n_p$ is not a valid constraint.

The model predicts a record $\hat{\mathbf{y}}=f_\Theta(\mathbf{H},\mathbf{c})$ and a verbalizer $V$ produces the description $\hat{s}=V(\hat{\mathbf{y}})$.
Our design favors traceability: language quality is determined by the correctness and completeness of $\hat{\mathbf{y}}$, not by surface-form similarity alone.

\subsection{Configuration-aware CSI encoder}

We standardize complex CSI features per example and project each beam's frequency response into a token.
The input projection combines a linear frequency projection, a depthwise--pointwise convolutional path, and adaptive average/max summaries.
Configuration metadata modulates the features with a feature-wise affine transformation.
Beam positional embeddings and encodings of carrier/frequency settings are then added.

A six-layer Transformer encoder \citep{vaswani2017attention} with width 384, six attention heads, and a 1536-dimensional feed-forward block processes the sequence together with a learned \texttt{[CLS]} token.
The \texttt{[CLS]} representation and the masked token mean are concatenated and projected to a 256-dimensional normalized CSI embedding $\mathbf{z}_c$.

\subsection{Language and prototype alignment}

Descriptions are tokenized with a 300-token vocabulary and maximum length 48.
A three-layer, four-head Transformer of width 256 produces a normalized text embedding $\mathbf{z}_t$.
We construct semantic prototypes from coarse channel keys that group examples by environment, \los/\nlos\ state, path richness, delay and angle regimes, $K$-factor bin, and interaction statistics.
After removing classes with fewer than five training samples, eight prototypes remain.

For a minibatch, we optimize symmetric InfoNCE losses for CSI--text, CSI--prototype, and text--prototype pairs:
\begin{equation}
\mathcal{L}_{a,b} =
 -\frac{1}{2N}\sum_i
 \left[
 \log\frac{\exp(\mathbf{z}_a^i{}^\top\mathbf{z}_b^i/T)}
 {\sum_j\exp(\mathbf{z}_a^i{}^\top\mathbf{z}_b^j/T)}
 +
 \log\frac{\exp(\mathbf{z}_b^i{}^\top\mathbf{z}_a^i/T)}
 {\sum_j\exp(\mathbf{z}_b^i{}^\top\mathbf{z}_a^j/T)}
 \right],
\end{equation}
with initial temperature $T=0.07$.
A semantic classifier and a balanced $K$-factor attribute classifier further shape the representation.

\section{Physics-Aware Multi-Task Learning}
\label{sec:physics}

A residual multi-layer perceptron maps $\mathbf{z}_c$ to a shared physics token.
It feeds an auxiliary head for normalized global targets and a set of specialized branches.
Temporal targets use a delay-specific encoder and estimated power-delay-profile statistics.
Power prediction uses a separate power-feature encoder.
LoS and first-path angles use dedicated contexts and sine/cosine regression to respect angular periodicity.
First-path delay, delay spread, first-path power, and strong-$K$ regimes combine coarse bin classification with within-bin position regression.
Reflection interactions use classification and regression heads.
The path-count extension adds a distinct regressor for $n_{rp}$.

The total objective is
\begin{align}
\mathcal{L}={}&
\lambda_{ct}\mathcal{L}_{c,t}+
\lambda_{cp}\mathcal{L}_{c,p}+
\lambda_{tp}\mathcal{L}_{t,p}+
\lambda_{\mathrm{sem}}\mathcal{L}_{\mathrm{sem}}+
\lambda_{\mathrm{attr}}\mathcal{L}_{\mathrm{attr}} \nonumber\\
&+\sum_{q\in\mathcal{Q}}\lambda_q\mathcal{L}_q+
\lambda_{\mathrm{rel}}\mathcal{L}_{\mathrm{rel}},
\end{align}
where $\mathcal{Q}$ contains target-specific regression, bin-classification, and within-bin losses.
The three alignment losses use weight 1.0; semantic classification, attribute classification, and auxiliary regression use 0.5, 0.1, and 0.005.
The principal physics weights are 0.08 for first-path delay, 0.05 each for LoS delay and LoS angle, 0.02 for first-path angle, and 0.02/0.01 for reflection classification/regression.
The complete loss specification is provided in Appendix~\ref{app:weights}.

Relational training penalties impose only relations justified by the label definitions: nonnegative spreads and interaction counts, at least one propagation path, positive LoS delay when LoS is present, and $\tau_1\ge\tau_{\los}$ for LoS examples.
We deliberately do not impose $n_r\le n_p$.

\section{CSI-to-Text Verbalization}
\label{sec:verbalization}

The final system does not ask a large language model to invent a caption.
Instead, it first predicts typed attributes and then realizes them with an explicit template.
A typical output is:
\begin{quote}\small
\emph{This signal is indoor and NLoS, with 5 propagation paths. No direct LoS path is detected. The first-path delay is 82.4 ns, the first-path angle is $-31.7^\circ$, the $K$-factor is 3.2 dB, and 4 paths contain reflection interactions.}
\end{quote}
Conditional clauses prevent LoS-only values from being stated for \nlos\ samples.
This design is less linguistically diverse than open-ended generation, but it makes every statement attributable to one prediction and avoids evaluating harmless paraphrases with exact string match.

\section{Relational Correction}
\label{sec:correction}

Before verbalization, a bounds correction clamps impossible negative spreads and interaction counts and enforces $n_p\ge1$.
For LoS samples, a relational correction additionally replaces
\begin{equation}
\hat{\tau}_1 \leftarrow \max(\hat{\tau}_1,\hat{\tau}_{\mathrm{LoS}}),
\end{equation}
which encodes that the detected direct path cannot arrive after the earliest propagation path.
This operation changes only inconsistent records and leaves the neural predictions otherwise untouched.
It is therefore a lightweight semantic safeguard rather than a second learned model.

\section{Experimental Setup}
\label{sec:setup}

\paragraph{Data.}
The main dataset contains 100,000 ray-traced examples with UPA-$8{\times}8$ measurements and 128 frequency samples.
It is balanced before splitting: 50,000 LoS and 50,000 NLoS examples.
We use a fixed 80/20 train/test split.
After removing semantic classes with fewer than five examples and filtering delay spread $\ge400$ ns, 79,429 training and 19,877 test examples remain.
All main-model, baseline, and ablation comparisons use this same split.

The reflected-path extension is a separate dataset whose path labels were reconstructed and verified.
It contains 100,000 examples but is naturally imbalanced: 13,628 LoS and 86,372 NLoS.
Its 80,000/20,000 split contains 10,902/69,098 LoS/NLoS training examples and 2,726/17,274 test examples.
We report this extension separately and do not reinterpret it as the balanced main dataset.

\paragraph{Optimization and compute.}
Models are trained for 100 epochs with AdamW, batch size 128, initial learning rate $3{\times}10^{-4}$, weight decay $10^{-2}$, five epochs of linear warmup, and cosine decay to $10^{-5}$.
Checkpoints are saved every 20 epochs.
Unless noted otherwise, results are the mean and sample standard deviation over seeds 0, 1, and 2.
The current implementation has 31.87M trainable parameters.
Each run uses one NVIDIA GeForce RTX 4090 with 23.6 GB memory; four such GPUs were available for parallel experiments.
The software stack is PyTorch 2.5.1+cu121 and CUDA 12.1.
Automatic mixed precision is disabled.
A representative run takes 5 h 27 min 44 s.

\paragraph{Metrics.}
For physical regression we report MAE; for reflection interactions we additionally report exact and adjacent-bin accuracy.
Text evaluation parses predicted and target descriptions into typed slots.
Categorical fields are measured by accuracy and macro-F1.
Numeric slot F1 evaluates whether the required values are present, while hallucination is the fraction of predicted numeric slots unsupported by the reference schema.
Tolerance accuracy measures whether a present predicted value lies within a field-specific tolerance: 50 ns for delays, $15^\circ$ for angles, 5 dB for first-path power, 3 dB for $K$, and one count for path/interaction fields.
The reported \emph{composite description factuality} is the macro average of categorical accuracies, numeric-slot F1, numeric tolerance accuracy, and $1-$hallucination.
It is not sentence-level exact match.
Physical consistency is an auxiliary rule-based diagnostic applied to the predicted record.

\section{Baseline Comparison}
\label{sec:baselines}

We compare with a flattened-CSI MLP, a convolutional baseline, an inverse-FFT power-delay-profile (PDP/IFFT) MLP, a Transformer without specialized physics branches, and the same CSI encoder trained with independent single-task heads.
Table~\ref{tab:baselines} reports the common three-seed evaluation.

\begin{table*}[t]
\centering
\caption{Physical prediction on the balanced main split (mean $\pm$ sample standard deviation over seeds 0, 1, and 2). Lower is better for MAE; higher is better for accuracy. ``Refl.'' denotes total reflection interactions, not reflected paths. Best results are bold.}
\label{tab:baselines}
\resizebox{\textwidth}{!}{
\begin{tabular}{lccccccc}
\toprule
Method &
$\tau_1$ MAE (ns) $\downarrow$ &
$K$ MAE (dB) $\downarrow$ &
Refl. MAE $\downarrow$ &
Exact $\uparrow$ &
Adjacent $\uparrow$ &
$\theta_1$ MAE ($^\circ$) $\downarrow$ &
$p_1$ MAE (dB) $\downarrow$\\
\midrule
\model\ full multi-task
& \textbf{21.911$\pm$0.303}
& 1.983$\pm$0.263
& \textbf{0.983$\pm$0.005}
& \textbf{0.7061$\pm$0.0039}
& \textbf{0.9742$\pm$0.0007}
& \textbf{18.372$\pm$0.365}
& 4.958$\pm$0.269\\
CSI encoder + single-task head
& 51.164$\pm$6.248
& \textbf{1.521$\pm$0.090}
& 1.222$\pm$0.072
& 0.5523$\pm$0.0259
& 0.9594$\pm$0.0126
& 20.449$\pm$0.431
& \textbf{3.740$\pm$0.209}\\
PDP/IFFT + MLP
& 28.750$\pm$4.611
& 5.002$\pm$0.275
& 1.706$\pm$0.032
& 0.4420$\pm$0.0201
& 0.9055$\pm$0.0019
& 40.292$\pm$0.159
& 12.236$\pm$0.278\\
Flattened CSI + MLP
& 114.371$\pm$0.978
& 5.600$\pm$1.013
& 1.835$\pm$0.196
& 0.4237$\pm$0.0999
& 0.8930$\pm$0.0431
& 32.782$\pm$0.413
& 13.315$\pm$1.221\\
Transformer without branches
& 268.177$\pm$0.003
& 22.762$\pm$0.000
& 3.214$\pm$0.021
& 0.1808$\pm$0.0000
& 0.5359$\pm$0.0000
& 53.111$\pm$0.005
& 38.090$\pm$0.013\\
CNN baseline
& 1608.970$\pm$1202.800
& 56.869$\pm$34.742
& 119.488$\pm$185.779
& 0.1562$\pm$0.1463
& 0.3107$\pm$0.1804
& 102.231$\pm$42.561
& 14591.900$\pm$18581.100\\
\bottomrule
\end{tabular}}
\end{table*}

The full model is best on five of seven metrics, including the propagation quantities most dependent on coupled temporal and angular structure.
The single-task system remains stronger for $K$-factor and first-path power, showing that multi-task sharing is not uniformly beneficial.
Compared with PDP/IFFT+MLP, the full system reduces first-path delay MAE by 23.8\% while also improving angle and reflection predictions.

\section{Three-Seed Text Factuality Results}
\label{sec:textresults}

\begin{table}[t]
\centering
\caption{Description quality on the reflected-path-enabled evaluation (mean $\pm$ sample standard deviation over three seeds).}
\label{tab:text}
\begin{tabular}{lc}
\toprule
Metric & Result\\
\midrule
Composite description factuality $\uparrow$ & 0.9406 $\pm$ 0.0107\\
Environment accuracy $\uparrow$ & 1.0000 $\pm$ 0.0000\\
LoS/NLoS accuracy $\uparrow$ & 0.9976 $\pm$ 0.0011\\
Categorical macro-F1 $\uparrow$ & 0.9975 $\pm$ 0.0011\\
Numeric-slot F1 $\uparrow$ & 0.9997 $\pm$ 0.0001\\
Numeric tolerance accuracy $\uparrow$ & 0.7063 $\pm$ 0.0548\\
Numeric-slot hallucination $\downarrow$ & 0.0004 $\pm$ 0.0002\\
Physical consistency $\uparrow$ & 0.9987 $\pm$ 0.0003\\
\bottomrule
\end{tabular}
\end{table}

Table~\ref{tab:text} shows that the system almost always emits the required slots and rarely introduces an unsupported one.
Categorical facts are also stable.
Numerical tolerance accuracy is lower and more seed-sensitive: the individual results are 69.84\%, 65.58\%, and 76.46\%.
This does not mean that one quarter of all sentences is entirely wrong; rather, it measures the proportion of jointly present numerical fields that fall within their field-specific tolerances.

The largest numerical errors occur for first-path angle (MAE $23.52\pm1.02^\circ$; tolerance accuracy $54.11\pm0.32\%$), first-path power ($9.19\pm1.69$ dB; $37.80\pm7.02\%$), $K$-factor ($3.65\pm1.70$ dB; $53.84\pm32.81\%$), and angle spread ($25.44\pm8.96^\circ$; $35.37\pm18.28\%$).
By comparison, first-path delay reaches $96.64\pm0.24\%$ accuracy at 50 ns and LoS angle reaches $99.60\pm0.13\%$ accuracy at $15^\circ$.

\section{Reflected-Path Extension}
\label{sec:pathcount}

The base model predicts the total number of reflection \emph{interactions}.
To support statements such as ``four paths contain reflection interactions,'' we add a separate reflected-path-count head.
On the imbalanced extension split, the verbalized reflected-path value attains an MAE of $0.091\pm0.007$, RMSE of $0.234\pm0.008$, and accuracy within one path of $98.82\pm0.06\%$.
These are text-level metrics computed after record generation and verbalization; they should not be confused with head-level exact classification accuracy.
The strong result demonstrates that the two reflection concepts can be modeled separately without relying on the invalid assumption that total interactions cannot exceed total paths.

\section{Ablation and Error Analysis}
\label{sec:ablation}

\paragraph{Architectural ablations.}
Table~\ref{tab:ablation} compares the full model with removal of the shared physics token or delay-specific encoder.
Removing the delay encoder degrades first-path delay, delay spread, $K$-factor, power, and reflection prediction, confirming the value of target-aware temporal features.
The shared token has a mixed effect: it improves LoS angle and reflection exact accuracy, but the no-shared variant is slightly better on several regression metrics.
We therefore interpret it as a useful coupling mechanism, not a universally superior component.

\begin{table}[t]
\centering
\caption{Selected three-seed ablations. Lower is better except reflection exact accuracy.}
\label{tab:ablation}
\resizebox{\columnwidth}{!}{
\begin{tabular}{lrrrr}
\toprule
Model & $\tau_1$ (ns) & $\sigma_\tau$ (ns) & $K$ (dB) & Refl. exact\\
\midrule
Full & 21.91$\pm$0.30 & 11.17$\pm$0.11 & 1.98$\pm$0.26 & \textbf{0.7061$\pm$0.0039}\\
No shared physics token & \textbf{21.62$\pm$0.71} & \textbf{11.15$\pm$0.08} & \textbf{1.74$\pm$0.04} & 0.7013$\pm$0.0041\\
No delay-specific encoder & 69.02$\pm$1.46 & 29.64$\pm$0.84 & 2.54$\pm$0.12 & 0.6652$\pm$0.0057\\
\bottomrule
\end{tabular}}
\end{table}

\paragraph{Effect of relational correction.}
In a representative run, bounds-only correction leaves 1,426 cases in which the predicted first path precedes the predicted LoS path.
Adding the relational correction reduces this count to zero and raises physical consistency from 92.64\% to 99.86\%.
Composite factuality changes from 95.6068\% to 95.6051\%, and numerical tolerance accuracy from 78.4292\% to 78.4205\%.
Thus, the rule removes contradictions at negligible measured accuracy cost.
We treat this consistency score as a diagnostic, not a substitute for comparison with ground-truth slots.

\paragraph{Failure modes.}
The remaining difficulty is continuous-value estimation rather than slot omission or category recognition.
Angle spread and $K$-factor also have large cross-seed variance.
Moreover, the path-count extension is dominated by NLoS samples, which are harder because their multipath structure is less directly observable.
Future evaluation should report LoS/NLoS-stratified errors and calibration intervals in addition to aggregate MAE.

\section{Limitations}

First, the output language is template-based and therefore factual and auditable but not stylistically open-ended.
The work demonstrates CSI-to-natural-language verbalization, not unrestricted dialogue about wireless scenes.
Second, the experiments use ray-traced data from one principal UPA/frequency configuration; generalization to measured channels, unseen arrays, bandwidths, and environments remains to be established.
Third, multi-task learning does not dominate single-task learning for every variable, especially $K$-factor and first-path power.
Fourth, tolerance-based numerical accuracy depends on application-specific thresholds and should always be presented with raw error metrics.
Finally, the reflected-path extension is strongly NLoS-skewed, so its aggregate score should not be compared directly with balanced main-split results.

\section{Conclusion}

We presented a physics-grounded framework that converts CSI into an interpretable natural-language channel description.
The system combines configuration-aware CSI encoding, language and prototype alignment, specialized multi-task physics heads, explicit verbalization, and a lightweight relational correction.
Three-seed experiments show strong physical prediction, near-complete slot realization, very low hallucination, and high post-correction consistency.
The reflected-path extension further distinguishes reflection interactions from the number of reflected paths.
The results support structured language as a practical interface for inspecting learned wireless channel models, while identifying continuous-value robustness and cross-configuration validation as the next priorities.

\subsubsection*{Reproducibility statement}
The anonymous supplementary code includes preprocessing, training, evaluation, signal-description parsing, and metric aggregation scripts.
Appendix~\ref{app:weights} records the optimization and loss configuration, and Section~\ref{sec:setup} reports data filtering, seeds, compute, and software versions.

\bibliography{references}
\bibliographystyle{iclr2026_conference}

\appendix

\section{Loss weights}
\label{app:weights}

\begin{table*}[h]
\centering
\caption{Loss configuration used for the main full model.}
\label{tab:weights}
\begin{tabular}{lll}
\toprule
Group & Term & Weight\\
\midrule
Alignment & CSI--text / CSI--prototype / text--prototype & 1.0 / 1.0 / 1.0\\
Representation & semantic classifier / attribute classifier / auxiliary physics regression & 0.5 / 0.1 / 0.005\\
$K$-factor & strong-bin classifier / within-bin position & 0.5 / 0.1\\
First-path delay & normalized / raw / fused raw & 0.08 / 0.001 / 0.02\\
First-path delay & bin classifier / position / consistency / tail underestimate & 0.02 / 0.01 / 0.01 / 0.01\\
Delay spread & raw / bin classifier / within-bin position & 0.0005 / 0.03 / 0.008\\
LoS delay & regression / nonnegative / consistency & 0.05 / 0.005 / 0.005\\
Angles & LoS / first path / first-path NLoS & 0.05 / 0.02 / 0.01\\
First-path power & bin classifier / position; NLoS sample multiplier & 0.03 / 0.1; 2.0\\
Reflection interactions & classifier / regression; NLoS sample multiplier & 0.02 / 0.01; 1.0\\
\bottomrule
\end{tabular}
\end{table*}

\section{Field-specific text metrics}
\label{app:textmetrics}

\begin{table*}[h]
\centering
\caption{Three-seed numerical text metrics on the reflected-path-enabled evaluation. Accuracy uses the tolerance shown in parentheses.}
\label{tab:fieldmetrics}
\begin{tabular}{lcc}
\toprule
Field & MAE & Tolerance accuracy\\
\midrule
Path count ($\pm1$) & 0.326$\pm$0.056 & 0.8893$\pm$0.0485\\
First-path delay ($\pm50$ ns) & 13.923$\pm$0.578 ns & 0.9664$\pm$0.0024\\
First-path angle ($\pm15^\circ$) & 23.517$\pm$1.015$^\circ$ & 0.5411$\pm$0.0032\\
First-path power ($\pm5$ dB) & 9.186$\pm$1.693 dB & 0.3780$\pm$0.0702\\
$K$-factor ($\pm3$ dB) & 3.646$\pm$1.699 dB & 0.5384$\pm$0.3281\\
Delay spread ($\pm50$ ns) & 11.429$\pm$0.076 ns & 0.9741$\pm$0.0015\\
Angle spread ($\pm15^\circ$) & 25.437$\pm$8.960$^\circ$ & 0.3537$\pm$0.1828\\
LoS delay ($\pm50$ ns) & 10.954$\pm$0.245 ns & 0.9783$\pm$0.0013\\
LoS angle ($\pm15^\circ$) & 1.321$\pm$0.148$^\circ$ & 0.9960$\pm$0.0013\\
Reflection interactions ($\pm1$) & 0.931$\pm$0.017 & 0.6499$\pm$0.0103\\
Reflected paths ($\pm1$) & 0.091$\pm$0.007 & 0.9882$\pm$0.0006\\
\bottomrule
\end{tabular}
\end{table*}

\end{document}
